\begin{thebibliography}{99}

\bibitem{AlpogeFurman2026}
L. Alp\"oge and R. Furman,
\emph{More than two thirds of the zeta zeros are simple and on the critical line},
arXiv:2608.13637 (2026).

\bibitem{BCY2011}
H.~M. Bui, J.~B. Conrey, and M.~P. Young,
\emph{More than 41\% of the zeros of the zeta function are on the critical line},
Acta Arith. \textbf{150} (2011), 35--64.

\bibitem{Conrey1989}
J.~B. Conrey,
\emph{More than two fifths of the zeros of the Riemann zeta function are on the critical line},
J. Reine Angew. Math. \textbf{399} (1989), 1--26.

\bibitem{Feng2012}
S. Feng,
\emph{Zeros of the Riemann zeta function on the critical line},
J. Number Theory \textbf{132} (2012), 511--542.

\bibitem{Hardy1914}
G.~H. Hardy,
\emph{Sur les z\'eros de la fonction $\zeta(s)$ de Riemann},
C. R. Acad. Sci. Paris \textbf{158} (1914), 1012--1014.

\bibitem{HLZ}
W. Heap, J. Li, and J. Zhao,
\emph{Lower bounds for discrete negative moments of the Riemann zeta function},
Algebr. Number Theory \textbf{16} (2022), 1589--1625;
arXiv:2003.09368.

\bibitem{HLP}
C.~P. Hughes, S. Lugmayer, and A. Pearce-Crump,
\emph{The second moment of the Riemann zeta function at its local extrema},
arXiv:2411.05573 (2024).

\bibitem{Levinson1974}
N. Levinson,
\emph{At least one-third of zeros of Riemann's zeta-function are on $\sigma=1/2$},
Proc. Natl. Acad. Sci. USA \textbf{71} (1974), 1013--1015.

\bibitem{Titchmarsh}
E.~C. Titchmarsh,
\emph{The Theory of the Riemann Zeta-Function},
2nd ed., revised by D.~R. Heath-Brown, Clarendon Press, Oxford, 1986.

\end{thebibliography}
